\documentclass[11pt]{letter}
\usepackage[margin=1in]{geometry}
\usepackage{times}
\usepackage{hyperref}
\usepackage[english]{babel}

\signature{Brendan Edmonds \\ UQ Cyber, School of EECS \\ University of Queensland \\ St Lucia QLD 4067, Australia \\ b.edmonds@uq.edu.au}
\address{Brendan Edmonds \\ UQ Cyber, School of Electrical Engineering and Computer Science \\ University of Queensland \\ St Lucia QLD 4067, Australia}

\begin{document}

\begin{letter}{The Editor-in-Chief \\ \emph{Journal of Software: Evolution and Process} \\ John Wiley \& Sons, Ltd}

\opening{Dear Editor,}

We are pleased to submit our research article, titled \emph{``Do Specifications Help LLMs Write Better Unit Tests?''}, for consideration as a regular research paper in the \emph{Journal of Software: Evolution and Process}.

The article addresses a question that sits at the intersection of two strands of long-standing interest to the journal's readership: the use of formal specifications as a tool for software evolution, and the increasingly pressing question of how to discipline large-language-model code generation so that the resulting artefacts are trustworthy enough to enter a production codebase. Concretely, we ask whether providing an LLM with a specification of a method's intended behaviour measurably improves the unit tests it writes, with test quality assessed by mutation score on a corpus of 312 methods from Apache Commons Lang.

The answer is affirmative and the effect is large: specifications raise mutation score by 40.7~percentage points and test count by 272\% relative to signature-only generation ($p < 0.001$, Cohen's $d = 2.41$). To enable the comparison at scale, we developed JML-Inferrer, a heuristic static-analysis system that produces JML specifications organised by weakest-precondition and strongest-postcondition reasoning and validated through OpenJML's Extended Static Checker. The tool is not the contribution being evaluated---it is the means by which the empirical question becomes answerable for 312 methods rather than for a handful---and the article foregrounds the empirical result accordingly.

We believe this work is a fit for \emph{JSEP} for three reasons. First, the article directly addresses how a software process can absorb LLM-generated tests without sacrificing oracle quality, which we expect to be a topic of growing concern for industrial and research readers alike. Second, the inference instrument and the test-generation experiment together form a single pipeline, and the article reports both ends in detail---spec-accuracy validation against two independent raters and OpenJML discharge of every inferred clause---which is the depth of methodological reporting that the journal favours. Third, we provide a full replication package, including raw mutation logs, prompt templates, all 6{,}240 generated test suites, and the containerised toolchain, which we hope will support follow-up work in the journal's community.

This article has not been published previously and is not under consideration elsewhere. All authors have approved the manuscript and agree with its submission. We declare no conflicts of interest.

We thank you for your consideration and look forward to your response.

\closing{Yours sincerely,}

\end{letter}

\end{document}
